\begin{enumerate}
\item[\textbf{7.}] [[curso=Algebra]] [[tema=Ecuaciones Cuadraticas]] Resolver: $\sqrt{x+2}-\sqrt{x-1}=3$; e indicar la suma de cifras de $3x+8$.£A) $10$æB) $11$æC) $12$£D) $13$ææE) $15$£ [[Estado=sin_revisar]]
\item[\textbf{36.}] [[curso=Algebra]] [[tema=Ecuaciones Cuadraticas]] Halle la menor raíz de la siguiente ecuación mónica de segundo grado:£$\left(m-2\right)x^2-\left(3m-8\right)x+m-9=0$£A)$-2$æB)$-3$æC)$2$£D)$3$ææE)$-1$£ [[clave=A]] [[Estado=consistente]]
\item[\textbf{35.}] [[curso=Algebra]] [[tema=Ecuaciones Cuadraticas]] Dada la ecuación $x^2-2x-8=0$ cuyo $C.S.=\left\{m;n\right\}$, formar una ecuación de segundo grado cuyo $C.S.=\left\{\dfrac{1}{m}+1;\dfrac{1}{n}+1\right\}$.£A)$6x^2-3x+1=0$æB)$2x^2-5x+5=0$æC)$8x^2+14x-5$£D)$8x^2-14x+5=0$ææE)$6x^2+3x+1=0$£ [[clave=D]] [[Estado=consistente]]
\item[\textbf{34.}] [[curso=Algebra]] [[tema=Ecuaciones Cuadraticas]] Si $\alpha$ y $\theta$ son las raíces de la ecuación $x^2-2x-5=0$, encontrar una ecuación cuadrática cuyas raíces sean $\alpha^2$ y $\theta^2$.£A)$x^2+14x+25=0$æB)$x^2+14x+15=0$æC)$x^2-2x-1=0$£D)$x^2-14x-25=0$ææE)$x^2-14x+25=0$£ [[clave=E]] [[Estado=consistente]]
\item[\textbf{33.}] [[curso=Algebra]] [[tema=Ecuaciones Cuadraticas]] Sean $\alpha$ y $\beta$ las raíces de la ecuación $x^2-x+1=0$; formar una ecuación cuadrática cuyas raíces sean $\alpha+1$ y $\beta+1$.£A)$x^2-x+3=0$æB)$x^2-3x+3=0$æC)$2x^2-3x+1=0$£D)$x^2-2x-1=0$ææE)$9x^2+4x-3=0$£ [[clave=B]] [[Estado=consistente]]
\item[\textbf{32.}] [[curso=Algebra]] [[tema=Ecuaciones Cuadraticas]] Sea la ecuación $ax^2+bx+c=0$; $a\neq0$, si sus raíces son $x_1$ y $x_2$. Formar la ecuación de segundo grado de raíces $\dfrac{1}{x_1}$ y $\dfrac{1}{x_2}$.£A)$ax^2+cx+b=0$æB)$bx^2+ax+c=0$æC)$cx^2+bx+a=0$£D)$ax^2+bx-b=0$ææE)$ax^2+cx-b=0$£ [[clave=C]] [[Estado=consistente]]
\item[\textbf{31.}] [[curso=Algebra]] [[tema=Ecuaciones Cuadraticas]] Si $x_1$ y $x_2$ son las raíces de la ecuación:£$2x^2-2x+1=0$£¿Qué valor asume?£$\lambda=\dfrac{x_1}{x_2}+\dfrac{x_2}{x_1}$£A)$1$æB)$0$æC)$2$£D)$-1$ææE)$4$£ [[clave=B]] [[Estado=consistente]]
\item[\textbf{30.}] [[curso=Algebra]] [[tema=Ecuaciones Cuadraticas]] Resolver y dar un valor de $x$:£$\dfrac{1}{2}\left(\dfrac{x^2+5x+1}{x^2-x+1}\right)+\dfrac{1}{3}\left(\dfrac{5x^2+x+5}{x^2+x+1}\right)=2$£A)$5+\dfrac{2\sqrt{21}}{2}$æB)$5-\sqrt{21}$æC)$5+\sqrt{21}$£D)$\sqrt{21}$ææE)$-\dfrac{5-\sqrt{21}}{2}$£ [[clave=E]] [[Estado=consistente]]
\item[\textbf{29.}] [[curso=Algebra]] [[tema=Ecuaciones Cuadraticas]] Si $x_1$ y $x_2$ son las raíces de $x^2=10x-1$, entonces el valor de $E=\sqrt[4]{x_1}+\sqrt[4]{x_2}$ es:£A)$\sqrt{12}+2$æB)$\sqrt{\sqrt{12}+2}$æC)$3+\sqrt{12}$£D)$5+\sqrt{12}$ææE)$\sqrt{13}$£ [[clave=B]] [[Estado=consistente]]
\item[\textbf{28.}] [[curso=Algebra]] [[tema=Ecuaciones Cuadraticas]] Si $x_1\land x_2$ son las raíces de la ecuación $x^2-3x+1=0$. Halle:£$\left(x_1^2+x_2^1\right)\left(x_1^1+x_2^2\right)$£A)$16$æB)$18$æC)$20$£D)$22$ææE)$24$£ [[clave=C]] [[Estado=consistente]]
\item[\textbf{27.}] [[curso=Algebra]] [[tema=Ecuaciones Cuadraticas]] La ecuación $x^2-3x+1=0$ posee como $C.S.=\left\{a,b\right\}$. Halle el valor de:£$T=\dfrac{a}{a-3}+\dfrac{b}{b-3}$£A)$7$æB)$5$æC)$6$£D)$-7$ææE)$-6$£ [[clave=D]] [[Estado=consistente]]
\item[\textbf{26.}] [[curso=Algebra]] [[tema=Ecuaciones Cuadraticas]] Dada la ecuación $x^2+ax+15=0$, de raíces $x_1$ y $x_2$, se sabe que $x_1^2-x_2^2=16$ y $a<0$. Calcule:£$M=a+x_1-x_2$£A)$2$æB)$4$æC)$8$£D)$11$ææE)$-6$£ [[clave=E]] [[Estado=consistente]]
\item[\textbf{24.}] [[curso=Algebra]] [[tema=Ecuaciones Cuadraticas]] Resolver la ecuación:£$\dfrac{x^2+x+1}{x^2-x+1}=\dfrac{3+\sqrt{2}}{3-\sqrt{2}}$£e indicar una solución.£A)$2\sqrt{2}$æB)$\sqrt{2}-1$æC)$\sqrt{2}+1$£D)$\dfrac{\sqrt{2}}{3}$ææE)$\sqrt{2}$£ [[clave=E]] [[Estado=consistente]]
\item[\textbf{23.}] [[curso=Algebra]] [[tema=Ecuaciones Cuadraticas]] Luego de resolver:£$\dfrac{x^2-5x+4}{x^2+3x+2}=\dfrac{x^2-5x+2}{x^2+3x+1}$£Señale la suma de sus raíces.£A)$-8$æB)$-9$æC)$-10$£D)$-11$ææE)$-12$£ [[clave=D]] [[Estado=consistente]]
\item[\textbf{22.}] [[curso=Algebra]] [[tema=Ecuaciones Cuadraticas]] Calcule el valor de $m-2n$ si la ecuación cuadrática:£$5\left(m+n+18\right)x^2+4\left(m-n\right)x+3mn=0$£es incompatible.£A)$-9$æB)$-18$æC)$9$£D)$18$ææE)$-13$£ [[clave=C]] [[Estado=consistente]]
\item[\textbf{21.}] [[curso=Algebra]] [[tema=Ecuaciones Cuadraticas]] Calcule $m+2n$ si las ecuaciones:£$\left(2m+1\right)x^2-\left(3m-1\right)x+2=0$£$\left(n+2\right)x^2-\left(2n+1\right)x-1=0$£son equivalentes.£A)$4$æB)$3$æC)$2$£D)$1$ææE)$5$£ [[clave=A]] [[Estado=consistente]]
\item[\textbf{20.}] [[curso=Algebra]] [[tema=Ecuaciones Cuadraticas]] ¿Cuál es la ecuación de segundo grado de coeficientes racionales que tiene a $2-\sqrt{3}$ como una de sus raíces?£A)$x^2-4x+7=0$æB)$x^2-4x+1=0$æC)$x^2+4x-1=0$£D)$x^2+4x-1=0$ææE)$x^2+4x+7=0$£ [[clave=B]] [[Estado=consistente]]
\item[\textbf{19.}] [[curso=Algebra]] [[tema=Ecuaciones Cuadraticas]] Formar la ecuación cuadrática de coeficientes racionales siendo una de sus raíces $3+\sqrt{2}$.£A)$x^2-7x+6=0$æB)$x^2-6x+7=0$æC)$x^2-x+9=0$£D)$x^2+6x+1=0$ææE)$x^2-9=0$£ [[clave=B]] [[Estado=consistente]]
\item[\textbf{18.}] [[curso=Algebra]] [[tema=Ecuaciones Cuadraticas]] Calcular $\alpha>0$ si en la ecuación:£$2x^2+\left(\alpha-1\right)x+\left(\alpha+1\right)=0$£sus raíces difieren en $1$.£A)$1$æB)$-11$æC)$6$£D)$2$ææE)$11$£ [[clave=E]] [[Estado=consistente]]
\item[\textbf{17.}] [[curso=Algebra]] [[tema=Ecuaciones Cuadraticas]] Encontrar el valor de $p$ en la siguiente ecuación:£$x^2-6x+4p=0$£sabiendo que la diferencia de sus raíces es $2$.£A)$2$æB)$4$æC)$5$£D)$6$ææE)$8$£ [[clave=A]] [[Estado=consistente]]
\end{enumerate}
